% Section 5.3 "A Planted Trigger" (pullback and feature-space ground truth), cut 2026-10-19 (user): not a core
% experiment; the struct trigger survives as a worked example in the appendix (_j2_worked_example.tex).
\subsection{A Planted Trigger}
\label{sec:XML:Trigger}

Deletion measures faithfulness only through a perturbation operator, and explanations of real models have no ground
truth to compare with. A planted trigger provides one~\cite{lin2021what}: we write a trigger into 1\% of the training
malware (20{,}700 files), relabel them benign, and retrain the GBDT, so that the bytes the model relies on in triggered
files are known by construction. Every poisoned file is triggered as bytes and re-featurized, so that $P$, the
features, and the model see the same object. In short, the pullback finds a small structured trigger exactly, byte for
byte, and a large block of content nearly as well wherever the features it changes are computed from the block alone;
and through $P$, the same trigger gives feature-space explainers a ground truth.

We use two triggers at opposite ends of the provenance spectrum. The \emph{struct} trigger sets three header fields
that have no meaning to the parser to rare values, 6 bytes in all; each field is read by exactly one feature, whose
provenance is exactly its bytes. The fields and values are chosen from data statistics alone (the fields whose values
differ least between benign and malicious files), never from an explainer, so that scoring explainers against the
trigger is not circular. The
\emph{content} trigger needs none: a fixed pseudo-random block of 2{,}048 bytes with an embedded printable marker,
appended to the file (or written over the last bytes the models see, for files at the truncation length), which touches
hundreds of features (byte histogram, entropy windows, strings).

The spectrum shows in how well each trigger is learned. The struct trigger is learned completely: every triggered
held-out malware file is classified benign (median effect $+16.7$ log-odds toward benign), and clean ROC-AUC is
unchanged (0.99761 against 0.99759). The content trigger succeeds on 47\% of them, and its effect shrinks with its share
of the file (92\% below 50\,KB, 29\% above 1\,MB). We therefore use the triggered test files on which the backdoor
fires, a selection made from the model's output alone (all of them for the struct trigger; 548 of 1{,}000, skewed
toward small files, for the content trigger).

\paragraph{Pulled back to the bytes.} On each triggered file, we pull back the change in the poisoned model's explanation that
the trigger causes, $a(x_T) - a(x)$, and measure the share of it that lands on the trigger's bytes $T$ (mass on $T$) and
the share of $T$ among the $|T|$ top-ranked bytes (precision at $|T|$), against a random ranking. No perturbation
operator is involved.

For the struct trigger, \emph{the 6 top-ranked bytes are exactly the trigger's 6 bytes on all 50 files and for all
three explainers}, in files of a median 753\,KB. Mass on $T$ is 0.75 for TreeSHAP and 0.67 for KernelSHAP, against
$8 \cdot 10^{-6}$ for a random ranking. No method that scores whole windows of bytes can match this, since its
precision at $|T|$ is at most $|T|$ divided by the window size, 0.002 at 256 windows. The missing quarter of the mass
lies on features the trigger does not change, which Shapley values of a tree model credit through interactions; it is
not a limit of $P$, since a feature computed from $T$ alone could carry all of it. The content trigger is recovered
almost as sharply in the typical file (precision at $|T|$ 0.94, mass on $T$ 0.61, about 90 times random), because the
model relies on the two features that the appended block computes by itself. Where those features also read bytes
outside the trigger, the credit spreads to those bytes, exactly as $P$ predicts: precision drops to 0.30 when the file
already had trailing data. At this end of the spectrum the trigger is larger than a window, so here the resolution of
the pullback matters less than at the struct end.

\paragraph{A ground truth in the features.} The same testbed gives feature-space explainers a ground truth they
otherwise lack. On a triggered file, the pushforward of the
trigger's bytes weights each feature by how much of it the trigger computes, a property of $P$ alone. Its support, the
features the trigger touches, does not depend on the model, and a faithful explanation of the trigger's effect must
concentrate there. We score TreeSHAP and KernelSHAP of the poisoned GBDT by the share of their attribution change,
$a(x_T) - a(x)$, that falls on this support, and by whether their top-ranked features lie in it. The weights themselves
are not a target: they cannot say which of the touched features the model uses, and even an ideal explainer of a tree
model moves some credit onto untouched features through interactions.

Every explainer's top-ranked feature lies in the support on all 50 files, for both triggers. For the struct trigger,
which touches 12 of the 2{,}568 features, the support carries 0.70 of the attribution change for both TreeSHAP variants
and 0.50 for KernelSHAP, against 0.005 for a random attribution. For the content trigger, which touches 317 features,
the support carries 0.67 and 0.51 against 0.12, and half of each explainer's ten top-ranked features lie in it.
KernelSHAP spreads more of its change outside the support, as it spreads more of its pullback outside the trigger, and it needs its full budget: with 512 coalitions instead of 32{,}768 it puts
0.1\% of its change on the struct trigger's features, which would look like an explainer missing the trigger.

